**Example 1**

Let us consider the operator $A$:

$$
A =
\begin{pmatrix}
1 & 0 & 0 \\
0 & 1 & 0 \\
0 & 0 & 2
\end{pmatrix}.
$$

To verify that it is a Hermitian matrix, we can verify its conjugate transpose:

$$
A^\dagger =
\begin{pmatrix}
1 & 0 & 0 \\
0 & 1 & 0 \\
0 & 0 & 2
\end{pmatrix}
= A.
$$

Its eigenvalues are:

$$
[1,\;1,\;2].
$$

And the corresponding eigenvectors (each column is an eigenvector) are

$$
\begin{pmatrix}
1 & 0 & 0 \\
0 & 1 & 0 \\
0 & 0 & 1
\end{pmatrix}.
$$

We have then that eigenvalue $1$ is degenerate. We can reconstruct matrix $A$ by spectral decomposition:

$$
A =
\begin{pmatrix}
1 & 0 & 0 \\
0 & 1 & 0 \\
0 & 0 & 2
\end{pmatrix}.
$$

We have now defined a Hermitian matrix $B$:

$$
B =
\begin{pmatrix}
3 & 1 & 0 \\
1 & 4 & 0 \\
0 & 0 & 5
\end{pmatrix}.
$$

We can verify that they commute, that is, $AB-BA=0$:

$$
AB-BA =
\begin{pmatrix}
0 & 0 & 0 \\
0 & 0 & 0 \\
0 & 0 & 0
\end{pmatrix}.
$$

In the simplest case, we then have a subspace of only one dimension for eigenvalue $2$. We can see that $Bv=bv$, with $b=5$:

$$
B
\begin{pmatrix}
0\\
0\\
1
\end{pmatrix}
=
\begin{pmatrix}
0\\
0\\
5
\end{pmatrix}.
$$

Now for eigenvalue $1$, we have a two-dimensional subspace formed by the vectors

$$
\begin{pmatrix}
1 & 0\\
0 & 1\\
0 & 0
\end{pmatrix}.
$$

All non-null vectors in this subspace continue to be eigenvectors of $A$ with the same eigenvalue $1$.

If we denote the two eigenvectors by $(v_1,v_2)$ and use them as a basis, any non-null vector that is a linear combination of them,

$$
u = av_1+bv_2,
$$

is also an eigenvector of $A$. In fact,

$$
Au
=
A(av_1+bv_2)
=
aAv_1+bAv_2
=
av_1+bv_2
=
u.
$$

Now applying $B$ to the first eigenvector, we can see that it is a linear combination of the other two vectors in the subspace and is not a $B$ eigenvector:

$$
B
\begin{pmatrix}
1\\
0\\
0
\end{pmatrix}
=
\begin{pmatrix}
3\\
1\\
0
\end{pmatrix}.
$$

That is, even though these two eigenvectors associated with $1$ are a basis for this subspace, they do not form a common eigenvector basis. But we can build one.

Since all non-null vectors in this subspace are eigenvectors of $A$, we only need to find those who are also eigenvectors of $B$.

In this two-dimensional space, $B$'s action on the vectors of the space is given by:

$$
B_1 =
\begin{pmatrix}
3 & 1\\
1 & 4
\end{pmatrix}.
$$

Its eigenvalues are

$$
[2.38,\;4.62].
$$

And the corresponding eigenvectors (each column is an eigenvector) are:

$$
\begin{pmatrix}
-0.85 & -0.53\\
0.53 & -0.85
\end{pmatrix}.
$$

Combining with the previous results, we can have a new basis given by the eigenvectors (where each column corresponds to a vector of the new basis):

$$
\begin{pmatrix}
-0.85 & -0.53 & 0\\
0.53 & -0.85 & 0\\
0 & 0 & 1
\end{pmatrix}.
$$

We can test by applying $A$ to the vectors of the new basis:

$$
A
\begin{pmatrix}
-0.85\\
0.53\\
0
\end{pmatrix}
=
\begin{pmatrix}
-0.85\\
0.53\\
0
\end{pmatrix},
$$

$$
A
\begin{pmatrix}
-0.53\\
-0.85\\
0
\end{pmatrix}
=
\begin{pmatrix}
-0.53\\
-0.85\\
0
\end{pmatrix},
$$

$$
A
\begin{pmatrix}
0\\
0\\
1
\end{pmatrix}
=
2 \begin{pmatrix}
0\\
0\\
1
\end{pmatrix}.
$$

With the corresponding eigenvalues:

$$
[1,\;1,\;2].
$$

We can test by applying $B$ to the vectors of the new basis:

$$
B
\begin{pmatrix}
-0.85\\
0.53\\
0
\end{pmatrix}
=
2.38
\begin{pmatrix}
-0.85\\
0.53\\
0
\end{pmatrix},
$$

$$
B
\begin{pmatrix}
-0.53\\
-0.85\\
0
\end{pmatrix}
=
4.62
\begin{pmatrix}
-0.53\\
-0.85\\
0
\end{pmatrix},
$$

$$
B
\begin{pmatrix}
0\\
0\\
1
\end{pmatrix}
=
5
\begin{pmatrix}
0\\
0\\
1
\end{pmatrix}.
$$

With the corresponding eigenvalues:

$$
[2.38,\;4.62,\;5].
$$


**Example 2**

Suppose we have the operators

$$
A=
\begin{pmatrix}
0 & 1\\
1 & 0
\end{pmatrix}
$$

and

$$
B=
\begin{pmatrix}
1 & 0\\
0 & -1
\end{pmatrix}.
$$

We can verify that they do not commute, since

$$
[A,B]=AB-BA
=
\begin{pmatrix}
0 & -2\\
2 & 0
\end{pmatrix}
\neq 0.
$$

For the operator \(A\), the eigenvalues are

$$
a=+1,\,-1,
$$

with corresponding normalized eigenvectors

$$
|a=+1\rangle
=
\frac{1}{\sqrt{2}}
\begin{pmatrix}
1\\
1
\end{pmatrix},
\qquad
|a=-1\rangle
=
\frac{1}{\sqrt{2}}
\begin{pmatrix}
-1\\
1
\end{pmatrix}.
$$

Suppose, for example, that after measuring \(A\), we obtain the value \(a=+1\). The system is then in the corresponding eigenstate,

$$
|\psi\rangle
=
|a=+1\rangle
=
\frac{1}{\sqrt{2}}
\begin{pmatrix}
1\\
1
\end{pmatrix}.
$$

Now consider measuring \(B\). The eigenvalues of \(B\) are also

$$
b=+1,\,-1,
$$

with corresponding eigenvectors

$$
|b=+1\rangle
=
\begin{pmatrix}
1\\
0
\end{pmatrix},
\qquad
|b=-1\rangle
=
\begin{pmatrix}
0\\
1
\end{pmatrix}.
$$

The probability of obtaining a particular eigenvalue \(b\) is given by the squared magnitude of the inner product between the corresponding eigenvector and the state of the system:

$$
P(b)=|\langle b|\psi\rangle|^2.
$$

Therefore,

$$
P(b=+1)
=
|\langle b=+1|a=+1\rangle|^2
=
\frac{1}{2},
$$

and

$$
P(b=-1)
=
|\langle b=-1|a=+1\rangle|^2
=
\frac{1}{2}.
$$

Thus, there is a \(50\%\) probability of observing \(+1\) and a \(50\%\) probability of observing \(-1\).

In other words, since \(A\) and \(B\) do not commute and, in this example, do not share any eigenvectors, after measuring one of the observables, the outcome of a subsequent measurement of the other is intrinsically indeterminate within this mathematical formalism.

